% CCSS 7.SP.C.8 — Sample Spaces for Compound Events
\section{Sample Spaces for Compound Events}

\begin{learningGoals}
\begin{itemize}
    \item Represent sample spaces using lists, tables, and tree diagrams
    \item Identify all possible outcomes for compound events
\end{itemize}
\end{learningGoals}

\begin{conceptBox}{Sample Space}
A sample space lists all possible outcomes of an experiment.
\end{conceptBox}

\begin{workedExample}{Flipping Two Coins}
List all possible outcomes.

HH, HT, TH, TT
\answerHighlight{HH, HT, TH, TT}
\end{workedExample}

\tipBox{Tree diagrams help organize outcomes for compound events.}

\begin{practiceBox}{Sample Spaces}
\resetProblems

\prob List all outcomes for rolling a die and flipping a coin.
\answerExplain{$6 \times 2 = 12$ outcomes}{Each die roll ($1$–$6$) with heads or tails: $1H, 1T, 2H, 2T, ..., 6T$.}

\prob List all outcomes for picking a red or blue marble, then flipping a coin.
\answerExplain{$4$ outcomes}{Red-Head, Red-Tail, Blue-Head, Blue-Tail.}

\prob How many outcomes for flipping $3$ coins?
\answerExplain{$8$ outcomes}{$2^3 = 8$; HHH, HHT, HTH, HTT, THH, THT, TTH, TTT.}

\prob List all outcomes for rolling two dice.
\answerExplain{$36$ outcomes}{$6 \times 6 = 36$; each die has $6$ options.}

\prob Why use a tree diagram?
\answerExplain{To organize outcomes}{Tree diagrams show all possible paths for compound events.}

\end{practiceBox}
